\chapter{Comparaison des fonctions}\label{ch:b2:comparison}

L'analyse asymptotique --- l'art de remplacer une quantité compliquée
par une quantité simple augmentée d'une erreur contrôlée --- a été
amorcée dans le volume de première année avec les développements de
Taylor. Ce chapitre en fait une discipline à part entière~:
développements le long d'échelles générales, la comparaison
série--intégrale dans toute sa force asymptotique, la formule de
Stirling (démontrée complètement), et l'étude systématique des suites
définies implicitement. Ces techniques sont le pain quotidien de
l'analyse asymptotique, et tout chapitre ultérieur qui estime quoi que
ce soit --- séries, intégrales, probabilités --- se nourrit à cette
table.

\section{Relations de comparaison et échelles}

\begin{definition}\label{def:b2:comparison:landau}
Au voisinage d'un point $a$ ($a \in \R$ ou $\pm\infty$), pour des
fonctions (ou des suites, avec $n \to \infty$)~: $f = o(g)$, $f =
O(g)$, $f \sim g$ comme dans le volume de première année. Une
\emph{échelle de comparaison}\index{échelle de comparaison} en $a$ est
une famille de fonctions positives, comparables deux à deux, totalement
ordonnée par $o(\cdot)$ --- l'échelle usuelle en $+\infty$ étant
\[
x^{\alpha} (\ln x)^{\beta}
\qquad (\alpha, \beta \in \R),
\]
ordonnée lexicographiquement en $(\alpha, \beta)$, raffinée au besoin
par des exponentielles $\eu^{\gamma x}$.
\end{definition}

\begin{definition}[Développement asymptotique]\label{def:b2:comparison:expansion}
$f$ admet le \emph{développement asymptotique}\index{développement
asymptotique}
\[
f = c_1 \varphi_1 + c_2\varphi_2 + \dots + c_k \varphi_k +
o(\varphi_k)
\qquad (\varphi_{i+1} = o(\varphi_i) \text{ dans l'échelle})
\]
lorsque les restes successifs vérifient les estimations affichées. Les
coefficients sont alors uniques~: $c_1 = \lim f/\varphi_1$, et par
récurrence $c_{i+1} = \lim\,(f - \sum_{j \leq i}
c_j\varphi_j)/\varphi_{i+1}$.
\end{definition}

\begin{example}
Les développements de Taylor sont des \omterm{def:b2:comparison:expansion}{développements asymptotiques} le
long de l'échelle $(x - a)^k$ en $a$. Mais la notion est strictement
plus large~: en $+\infty$,
\[
\frac{1}{x - \ln x}
= \frac1x \cdot \frac{1}{1 - \frac{\ln x}{x}}
= \frac1x + \frac{\ln x}{x^2} + o\Bigl(\frac{\ln x}{x^2}\Bigr),
\]
un développement le long de l'échelle mixte --- aucun théorème de
Taylor ne s'applique, seulement le développement géométrique et le
calcul des $o$.
\end{example}

\begin{example}[L'échelle usuelle est réellement ordonnée]\label{ex:b2:comparison:scaleorder}
L'affirmation lexicographique de la~\cref{def:b2:comparison:landau}
demande une ligne de démonstration par cas. Comparons
$x^{\alpha}(\ln x)^{\beta}$ et $x^{\alpha'}(\ln x)^{\beta'}$
en $+\infty$. Si $\alpha < \alpha'$~: le rapport vaut
$x^{\alpha - \alpha'}(\ln x)^{\beta - \beta'} \to 0$, car une
puissance négative de $x$ écrase toute puissance de $\ln x$ (poser $x =
\eu^t$~: $\eu^{(\alpha - \alpha')t}\,t^{\beta - \beta'} \to 0$
par la limite «~l'exponentielle l'emporte sur le polynôme~» du volume
de première année). Si $\alpha = \alpha'$ et $\beta < \beta'$~: le
rapport vaut $(\ln x)^{\beta - \beta'} \to 0$ directement. Ainsi les
couples $(\alpha, \beta)$, ordonnés lexicographiquement, ordonnent
l'échelle par $o(\cdot)$ --- et la substitution $x = \eu^t$ est
l'astuce universelle pour les comparaisons mixtes
puissance-logarithme.
\end{example}

\begin{example}[Classer une ménagerie]\label{ex:b2:comparison:ranking}
Les échelles doivent être \emph{ordonnées}~; voici l'exercice type. En
$+\infty$, comparons $n^{10}$, $\eu^{\sqrt{\ln n}\,\cdot\,\sqrt n}$,
$2^n$ et $n^{\ln n}$ en prenant les logarithmes~:
\[
10\ln n
\;\ll\; (\ln n)^2
\;\ll\; \sqrt{n\ln n}
\;\ll\; n\ln 2 ,
\]
où $a_n \ll b_n$ signifie $a_n = o(b_n)$~; la deuxième entrée est
$\ln(n^{\ln n})$. Les exponentielles préservent ces écarts stricts (si
$\ln u_n - \ln v_n \to -\infty$ alors $u_n/v_n \to 0$), donc
\[
n^{10} = o\bigl(n^{\ln n}\bigr),
\qquad
n^{\ln n} = o\bigl(\eu^{\sqrt{n\ln n}}\bigr),
\qquad
\eu^{\sqrt{n\ln n}} = o(2^n) .
\]
La morale, doublement~: comparer toujours via les logarithmes
(différences de logarithmes, non rapports de logarithmes), et ne
jamais conclure $u_n \sim v_n$ à partir de $\ln u_n \sim \ln v_n$ ---
le couple $n^{10}$ et $n^{\ln n}$ a un rapport de $\ln$ tendant vers
$\infty$, mais $2^n$ et $4^n$ ont un rapport de $\ln$ exactement égal à
$2$ et sont follement non équivalents.
\end{example}

\section{Comparaison série--intégrale, asymptotiquement}

\begin{theorem}\label{thm:b2:comparison:seriesintegral}
Soit $f$ \omterm{def:b2:metric:continuity}{continue}, positive, \emph{décroissante} sur
$\intco{1}{+\infty}$.
\begin{enumerate}
  \item Si $\int_1^{\infty} f$ converge, les restes vérifient
        \[
        \int_{n+1}^{\infty} f \;\leq\; \sum_{k > n} f(k) \;\leq\;
        \int_{n}^{\infty} f .
        \]
  \item Si $\int_1^\infty f$ diverge, les sommes partielles vérifient
        $\sum_{k=1}^{n} f(k) = \int_1^n f + C + o(1)$ pour une
        certaine constante $C$~: la différence $\sum_{k \leq n} f(k) -
        \int_1^n f$ \emph{converge}.
\end{enumerate}
\end{theorem}

\begin{proof}
L'encadrement $f(k+1) \leq \int_k^{k+1} f \leq f(k)$ (décroissance)
était le procédé de première année~; en sommant sur $k \geq n+1$,
resp.\ $k \geq n$, on obtient (1). Pour (2), posons $u_k = f(k) -
\int_k^{k+1} f$~: par l'encadrement, $0 \leq u_k \leq f(k) - f(k+1)$,
donc les sommes partielles de $\sum u_k$ sont majorées par le
télescopage $f(1) - f(n+1) \leq f(1)$~: la série converge. De plus la
suite $\bigl(\int_n^{n+1} f\bigr)_n$ est décroissante ($f$ décroît) et
positive, donc convergente. En écrivant
\[
\sum_{k=1}^{n} f(k) - \int_1^n f
= \sum_{k=1}^{n} u_k + \int_n^{n+1} f ,
\]
le membre de droite converge quand $n \to \infty$~: la différence
converge vers une constante $C$, ce qui est l'énoncé (2).
\end{proof}

\begin{example}[Le développement harmonique]\label{ex:b2:comparison:harmonic}
Pour $f(t) = \frac1t$~: $H_n = \ln n + \gamma + o(1)$, ce qui retrouve
la constante d'Euler (volume de première année) avec une démonstration
plus propre. En poussant d'un \omterm{def:b2:structures:generated}{ordre} de plus
(\cref{exo:b2:comparison:3})~:
\[
H_n = \ln n + \gamma + \frac{1}{2n} + o\Bigl(\frac1n\Bigr).
\]
Les chiffres rendent le gain visible en $n = 10$~: $H_{10} =
2.928968\dots$ et $\ln 10 = 2.302585\dots$, donc l'estimation brute
de $\gamma$ est $H_{10} - \ln 10 = 0.626383$, à $0.049$ près~; en
soustrayant la correction $\frac1{20}$ on obtient $0.576383$, qui ne
diffère de $\gamma = 0.577216$ que de $8.3\cdot10^{-4}$ --- lequel est
lui-même le terme suivant $\frac{1}{12\cdot100}$ du développement,
comme le démontre le problème du week-end (question 8).
\end{example}

\begin{example}[Un $\ln(n!)$ grossier sans Stirling]\label{ex:b2:comparison:lnfactorial}
Le seul procédé d'encadrement localise déjà $\ln(n!)$. Comme
$\ln$ croît,
\[
\int_{k-1}^{k}\ln t\,\dd t \;\leq\; \ln k \;\leq\;
\int_{k}^{k+1}\ln t\,\dd t ,
\]
et en sommant sur $k = 2, \dots, n$ (avec $\int_1^n\ln = n\ln n
- n + 1$)~:
\[
n\ln n - n + 1 \;\leq\; \ln(n!) \;\leq\; (n+1)\ln(n+1) - n .
\]
Les deux bornes valent $n\ln n - n + O(\ln n)$~: d'où $\ln(n!) = n\ln
n - n + O(\ln n)$, et en particulier $\ln(n!) \sim n\ln n$.
Ce que Stirling ajoute, ce sont les deux échelons suivants --- le
$\frac12\ln n$ et la constante $\ln\sqrt{2\pi}$ --- qui coûtent le
télescopage plus fin du~\cref{thm:b2:comparison:stirling}. Savoir
quelle précision chaque outil achète est la moitié du métier de
l'asymptoticien.
\end{example}

\begin{example}[Les compensations exigent des développements]\label{ex:b2:comparison:cancellation}
Calculons la limite de $\sqrt{n^2 + n} - n$. Les deux termes sont
$\sim n$, et «~$\sim n - n$~» n'a aucun sens~: on ne peut pas
soustraire des équivalents. Développons plutôt~:
\[
\sqrt{n^2 + n} - n
= n\Bigl(\sqrt{1 + \tfrac1n} - 1\Bigr)
= n\Bigl(\frac{1}{2n} - \frac{1}{8n^2} +
O\Bigl(\frac1{n^3}\Bigr)\Bigr)
= \frac12 - \frac{1}{8n} + O\Bigl(\frac1{n^2}\Bigr) :
\]
limite $\frac12$, avec la vitesse d'approche $\frac1{8n}$ en prime. Le
mécanisme mérite un nom~: une différence de deux grandes quantités
\omterm{def:b2:nvs:equivalent}{équivalentes} vit entièrement dans leurs termes \emph{suivants}, si
bien qu'il faut développer jusqu'au premier \omterm{def:b2:structures:generated}{ordre} où les deux membres
diffèrent --- et traîner le reste pour certifier que rien d'autre ne
survit à cet \omterm{def:b2:structures:generated}{ordre}.
\end{example}

\begin{example}[Une comparaison divergente, en détail]\label{ex:b2:comparison:loglog}
Pour $f(t) = \frac{1}{t\ln t}$ sur $\intco{2}{+\infty}$
(\omterm{def:b2:metric:continuity}{continue}, positive, décroissante)~: $\int_2^x f = \ln\ln x -
\ln\ln 2 \to \infty$, donc d'après le~%
\cref{thm:b2:comparison:seriesintegral}~(2),
\[
\sum_{k=2}^{n}\frac{1}{k\ln k} = \ln\ln n + C + o(1)
\]
pour une certaine constante $C$. Deux leçons. Premièrement, la
divergence est réelle mais glaciale~: la somme partielle ne dépasse
$4$ que vers $n \approx \eu^{\eu^{4 - C}}$, astronomiquement grand.
Deuxièmement, la \emph{forme} $\ln\ln n$ a été livrée par une
primitive, non devinée~: pour des termes monotones, l'intégrale est le
procédé de sommation canonique, et la constante $C$ --- comme le
$\gamma$ d'Euler --- est la mémoire des premiers termes.
\end{example}

\section{Formule de Stirling}

\begin{lemma}[Intégrales de Wallis, revisitées]\label{lem:b2:comparison:wallis}
Soit $W_n = \int_0^{\pi/2} \sin^n t\,\dd t$. Alors $nW_nW_{n-1} =
\frac\pi2$ pour $n \geq 1$, $(W_n)$ décroît, et $W_n \sim
\sqrt{\dfrac{\pi}{2n}}$.
\end{lemma}

\begin{proof}
L'intégration par parties donne $nW_n = (n-1)W_{n-2}$ ($n \geq 2$),
donc $nW_nW_{n-1}$ est constant en $n$, égal à $1 \cdot W_1 W_0 =
\frac\pi2$. Décroissance~: $\sin^{n+1} \leq \sin^n$ sur
$\intcc{0}{\frac\pi2}$. L'encadrement, en détail~: la monotonie
donne $W_{n+1} \leq W_n \leq W_{n-1}$, et en divisant par $W_{n-1}
> 0$,
\[
\frac{n}{n+1} = \frac{W_{n+1}}{W_{n-1}} \leq
\frac{W_n}{W_{n-1}} \leq 1 ,
\]
l'identité de gauche venant de la récurrence à l'indice $n + 1$. Les
deux bornes tendent vers $1$~: $W_n \sim W_{n-1}$, d'où
\[
nW_n^2 \sim nW_nW_{n-1} = \frac\pi2
\qquad\Longrightarrow\qquad
W_n \sim \sqrt{\frac{\pi}{2n}} .
\]
\end{proof}

\begin{example}[Les premières intégrales de Wallis]\label{ex:b2:comparison:wallisvalues}
À partir de $W_0 = \frac\pi2$, $W_1 = 1$ et de la récurrence $nW_n =
(n-1)W_{n-2}$~:
\[
W_2 = \frac\pi4, \qquad
W_3 = \frac23, \qquad
W_4 = \frac{3\pi}{16}, \qquad
W_5 = \frac{8}{15}, \qquad
W_6 = \frac{5\pi}{32}.
\]
Les indices pairs portent un $\pi$, les impairs sont rationnels --- les
deux produits entrelacés des formes closes. Numériquement $W_6
\approx 0.4909$ contre l'asymptotique
$\sqrt{\pi/12} \approx 0.5116$~: en $n = 6$ l'équivalent est
déjà à $5\%$ près, et l'identité de produit est exacte pour tout
$n$~: $6\,W_6W_5 = 6\cdot\frac{5\pi}{32}\cdot\frac8{15} =
\frac\pi2$. De petites tables comme celle-ci sont le moyen le moins
coûteux d'attraper une erreur de calcul avant qu'elle n'infecte un
argument asymptotique.
\end{example}

\begin{theorem}[Stirling]\label{thm:b2:comparison:stirling}
\[
n! \;\sim\; \sqrt{2\pi n}\, \Bigl(\frac{n}{\eu}\Bigr)^{\!n}
.\index{formule de Stirling}
\]
\end{theorem}

\begin{proof}
\emph{Étape 1~: $n! \sim C \sqrt n\, (n/\eu)^n$ pour une certaine
constante $C > 0$.} Posons
\[
d_n = \ln(n!) - \Bigl(n + \frac12\Bigr)\ln n + n .
\]
Alors
\[
d_n - d_{n+1}
= \Bigl(n + \frac12\Bigr) \ln\frac{n+1}{n} - 1
= \Bigl(n + \frac12\Bigr)\Bigl(\frac1n - \frac{1}{2n^2} +
\frac{1}{3n^3} + o\bigl(n^{-3}\bigr)\Bigr) - 1
= \frac{1}{12n^2} + o\Bigl(\frac{1}{n^2}\Bigr),
\]
par le développement de Taylor de $\ln(1 + \frac1n)$. La série $\sum
(d_n - d_{n+1})$ converge donc absolument (comparaison avec $\sum
n^{-2}$), donc $(d_n)$ converge, disons vers $d$~; en exponentiant,
$n! \sim C\sqrt n\,(n/\eu)^n$ avec $C = \eu^{d}$.

\emph{Étape 2~: $C = \sqrt{2\pi}$ via Wallis.} La forme close $W_{2p}
= \frac{(2p)!}{4^p (p!)^2}\cdot\frac\pi2$ (issue de la récurrence,
calcul de première année refait dans le cadre du~%
\cref{lem:b2:comparison:wallis}) se combine avec l'étape 1~:
\[
W_{2p} \sim \frac{C\sqrt{2p}\,(2p/\eu)^{2p}}
{4^p\,\bigl(C\sqrt p\,(p/\eu)^p\bigr)^2}\cdot\frac{\pi}{2}
= \frac{\sqrt{2p}}{C\,p}\cdot\frac{\pi}{2}
= \frac{\pi}{C}\cdot\frac{1}{\sqrt{2p}} .
\]
En comparant avec $W_{2p} \sim \sqrt{\frac{\pi}{4p}}$
(\cref{lem:b2:comparison:wallis})~: $\frac{\pi}{C\sqrt{2p}} =
\sqrt{\frac{\pi}{4p}}\,(1 + o(1))$ force $C = \pi
\sqrt{\frac{4p}{2p\,\pi}} = \sqrt{2\pi}$.
\end{proof}

\begin{example}[Coefficient binomial central]\label{ex:b2:comparison:centralbinomial}
\[
\binom{2n}{n} = \frac{(2n)!}{(n!)^2}
\sim \frac{\sqrt{4\pi n}\,(2n/\eu)^{2n}}{2\pi n\,(n/\eu)^{2n}}
= \frac{4^n}{\sqrt{\pi n}} :
\]
la probabilité qu'une marche aléatoire symétrique revienne en $0$ à
l'instant $2n$ est $\sim \frac{1}{\sqrt{\pi n}}$ --- une annonce du~%
\cref{ch:b2:randomvar}.
\end{example}

\begin{remark}[Perspectives au sein de ce volume]\label{rem:b2:comparison:perspectives}
Tout chapitre quantitatif à venir parle la langue de ce chapitre.
Le~\cref{ch:b2:series} classe les séries en comparant les termes à
l'échelle $n^{-\alpha}(\ln n)^{-\beta}$ --- son problème du week-end
cartographie complètement cette frontière.
Le~\cref{ch:b2:integration} fait de même pour les intégrales
impropres, avec l'échelle identique dans la variable \omterm{def:b2:metric:continuity}{continue}.
Le~\cref{ch:b2:powerseries} calcule les rayons de convergence à partir
de $\limsup\abs{a_n}^{1/n}$, un exercice d'équivalents de racines
$n$-ièmes où Stirling est la clé standard ($\sqrt[n]{n!}
\sim \frac n\eu$, \cref{exo:b2:comparison:4}). Et les chapitres de
probabilités encaissent Stirling directement~: les estimations
locales du~\cref{ch:b2:randomvar} pour les coefficients binomiaux
sont exactement l'\cref{ex:b2:comparison:centralbinomial} et
l'\cref{ex:b2:comparison:lopsided}. L'asymptotique n'est pas
un chapitre ici~; c'est l'accent du volume.
\end{remark}

\begin{method}[La check-list de l'amorçage]\label{met:b2:comparison:checklist}
Avant de faire confiance à un développement obtenu par amorçage,
vérifiez quatre points.
(1) \emph{L'existence d'abord~:} la racine ou la suite doit être fixée
(monotonie, valeurs intermédiaires) avant tout développement
--- des symboles sans référent se développent magnifiquement et ne
signifient rien. (2) \emph{Un \omterm{def:b2:structures:generated}{ordre} par passe~:} chaque substitution
ne peut être crue qu'à l'\omterm{def:b2:structures:generated}{ordre} de l'estimation injectée~; extraire
deux termes nouveaux d'une seule passe est la source classique de
coefficients faux. (3) \emph{Les restes voyagent avec~:} traînez le
$o(\cdot)$ à travers chaque étape algébrique et laissez l'absorption
(les termes plus petits engloutis par des restes plus grands) se
produire à la fin, explicitement. (4) \emph{Audit numérique~:} évaluez
en une valeur honnête de $n$~; une erreur de coefficient survit
étonnamment souvent à une nouvelle dérivation algébrique, et ne survit
presque jamais à l'arithmétique.
\end{method}

\begin{remark}[Pièges courants]\label{rem:b2:comparison:pitfalls}
(i) Les équivalents s'\emph{additionnent} mal~: de $u_n \sim n + \ln n$
et $v_n \sim -n$ on ne peut \emph{pas} conclure $u_n + v_n \sim \ln
n$~; les compensations exigent des développements avec restes
explicites, jamais des équivalents nus. (ii) Ne jamais exponentier une
équivalence~: $n + 1 \sim n$ mais $\eu^{n+1} \not\sim \eu^n$~; la
direction sûre est de prendre les logarithmes d'équivalents tendant
vers $+\infty$ (problème du week-end de ce chapitre, question 24).
(iii) Un \omterm{def:b2:comparison:expansion}{développement asymptotique} est attaché à une \emph{échelle}~:
écrire $f = \frac1x + o\bigl(\frac1{x^2}\bigr)$ affirme plus que $f =
\frac1x + o\bigl(\frac1x\bigr)$, et mélanger les deux invalide
l'\omterm{def:b2:structures:algebra}{algèbre} qui suit. (iv) Dans les amorçages, substituer le
développement courant \emph{entier}, reste compris --- laisser tomber
un $o(\cdot)$ en cours de passe produit des coefficients plausibles
mais faux. (v) La comparaison série--intégrale requiert la monotonie~:
pour des termes oscillants elle échoue purement et simplement (comparer
$\sum\frac{\sin k}k$, \cref{ch:b2:series}).
\end{remark}

\begin{example}[Stirling en chiffres]\label{ex:b2:comparison:stirlingnum}
En $n = 10$~: la formule donne $\sqrt{20\pi}\,(10/\eu)^{10}
\approx 3\,598\,696$ contre $10! = 3\,628\,800$~: erreur relative
$8.3\cdot10^{-3}$, remarquable pour un énoncé «~asymptotique~» en
$n = 10$. L'erreur a une structure --- le raffinement exact $n!
= \sqrt{2\pi n}\,(n/\eu)^n\bigl(1 + \frac1{12n} +
O(n^{-2})\bigr)$ --- dont la première correction $\frac1{120} \approx
8.3\cdot10^{-3}$ explique l'écart observé presque exactement. La
machinerie d'Euler--Maclaurin du problème du week-end est précisément
la source systématique de tels termes de correction.
\end{example}

\begin{remark}[Où ce chapitre est utilisé]
La comparaison asymptotique est la grammaire de tout ce qui est
quantitatif en aval~: les critères de convergence et le panorama de
Bertrand du~\cref{ch:b2:series}, les critères d'intégrabilité du~%
\cref{ch:b2:integration}, les calculs de rayon de convergence du~%
\cref{ch:b2:powerseries}, et les théorèmes limites du~%
\cref{ch:b2:randomvar} (où Stirling fait tourner les estimations de de
Moivre--Laplace). Le volume de troisième année industrialise l'unique
idée que nous démontrons ici à la main --- extraire le terme
principal, majorer le reste --- en la méthode de Laplace et la
convergence dominée.
\end{remark}

\begin{example}[Une intégrale comparée à elle-même~: $\int_2^x \frac{\dd t}{\ln t}$]\label{ex:b2:comparison:liworked}
La boîte à outils de comparaison marche aussi sur les intégrales. Soit
$F(x) = \int_2^x\frac{\dd t}{\ln t}$ (l'intégrande est \omterm{def:b2:metric:continuity}{continu} sur
$\intco2\infty$). Intégrons par parties~:
\[
F(x) = \Bigl[\frac{t}{\ln t}\Bigr]_2^x +
\int_2^x\frac{\dd t}{(\ln t)^2}
= \frac{x}{\ln x} + O\Bigl(\int_2^x\frac{\dd t}{(\ln
t)^2}\Bigr) + O(1),
\]
et l'intégrale de reste est $o\bigl(\frac{x}{\ln x}\bigr)$~:
découpons-la en $\sqrt x$, en majorant par
\[
\int_2^{\sqrt x}\frac{\dd t}{(\ln t)^2} \leq \sqrt x
\qquad\text{et}\qquad
\int_{\sqrt x}^{x}\frac{\dd t}{(\ln t)^2} \leq
\frac{x}{(\ln\sqrt x)^2} = \frac{4x}{(\ln x)^2} .
\]
D'où $F(x) \sim \frac{x}{\ln x}$. Les lecteurs qui ont rencontré le
théorème des nombres premiers dans le problème du week-end de ce
chapitre reconnaîtront $F$~: c'est le logarithme intégral, le meilleur
estimateur de $\pi(x)$, et le calcul montre qu'il coïncide avec
$\frac{x}{\ln x}$ au premier \omterm{def:b2:structures:generated}{ordre}.
\end{example}

\begin{example}[Stirling sur un binomial déséquilibré]\label{ex:b2:comparison:lopsided}
La même routine à trois factorielles que pour
l'\cref{ex:b2:comparison:centralbinomial} donne, pour
$\binom{3n}{n} = \frac{(3n)!}{n!\,(2n)!}$~:
\[
\binom{3n}{n} \sim
\frac{\sqrt{6\pi n}\,(3n/\eu)^{3n}}
{\sqrt{2\pi n}\,(n/\eu)^{n}\cdot\sqrt{4\pi n}\,(2n/\eu)^{2n}}
= \sqrt{\frac{3}{4\pi n}}\,
\Bigl(\frac{27}{4}\Bigr)^{\!n} .
\]
Le taux exponentiel $\frac{27}4 = \frac{3^3}{2^2}$ vaut
$\eu^{3n\,H(1/3)}$ dans la notation d'entropie de la théorie de
l'information~: les binomiaux déséquilibrés croissent strictement plus
lentement que le $4^n$ central par deux pas --- ici $(27/4)^{1/3}
\approx 1.89 < 2$ par pas. Toute asymptotique binomiale en combinatoire
et en probabilités (\cref{ch:b2:randomvar}) est ce seul calcul avec
des poids différents.
\end{example}

\section{Suites définies implicitement}

\begin{method}\label{met:b2:comparison:implicit}
Pour trouver l'asymptotique des solutions $x_n$ d'une équation $F(x, n) =
0$~:\index{suite implicite}
\begin{enumerate}
  \item \emph{Localiser}~: démontrer l'existence et l'unicité de $x_n$
        dans un intervalle précis (monotonie, théorème des valeurs
        intermédiaires), et trouver son comportement grossier (limite,
        \omterm{def:b2:structures:generated}{ordre} de grandeur).
  \item \emph{Amorcer}~: substituer la forme grossière $x_n = (\text{terme
        principal})(1 + \varepsilon_n)$ dans l'équation et résoudre
        pour l'\omterm{def:b2:structures:generated}{ordre} suivant de $\varepsilon_n$~; répéter, chaque passe
        raffinant d'un \omterm{def:b2:structures:generated}{ordre}.
\end{enumerate}
\end{method}

\begin{example}\label{ex:b2:comparison:tan}
Pour $n \geq 1$, l'équation $\tan x = x$ a exactement une solution
$x_n$ dans $\intoo{n\pi - \frac\pi2}{n\pi + \frac\pi2}$ (la fonction
$\tan x - x$ y croît de $-\infty$ à $+\infty$, sa dérivée valant
$\tan^2 x \geq 0$). \emph{Grossièrement~:} $x_n = n\pi +
\frac\pi2 - y_n$ avec $y_n \in \intoo{0}{\pi}$~; comme $x_n \to
\infty$ et $\tan x_n = x_n \to +\infty$, $x_n$ approche l'asymptote
par la gauche~: $y_n \to 0$. \emph{Amorçage~:} $\tan x_n =
\cot y_n = \frac{1}{\tan y_n} \sim \frac{1}{y_n}$, et l'équation
$\cot y_n = x_n \sim n\pi$ donne $y_n \sim \frac{1}{n\pi}$. D'où
\[
x_n = n\pi + \frac\pi2 - \frac{1}{n\pi} +
o\Bigl(\frac1n\Bigr),
\]
et le procédé se poursuit à tout \omterm{def:b2:structures:generated}{ordre}
(\cref{exo:b2:comparison:6}).
\end{example}

\begin{example}[Un second passage de la méthode]\label{ex:b2:comparison:xplusln}
Résolvons $x + \ln x = n$ asymptotiquement. \emph{Localiser~:} $x
\mapsto x + \ln x$ croît de $-\infty$ à $+\infty$ sur
$\intoo{0}{+\infty}$~: une racine unique $x_n$, et $x_n \to \infty$.
\emph{Grossièrement~:} $\ln x_n = o(x_n)$ donne $x_n \sim n$.
\emph{Amorçage~:} de $x_n = n - \ln x_n$ et $\ln x_n = \ln n
+ o(1)$ (logarithmes d'équivalents, les deux membres $\to \infty$)~:
\[
x_n = n - \ln n + o(1) ;
\]
une passe de plus, avec $\ln x_n = \ln\bigl(n - \ln n + o(1)\bigr) =
\ln n - \frac{\ln n}{n} + o\bigl(\frac{\ln n}n\bigr)$~:
\[
x_n = n - \ln n + \frac{\ln n}{n} +
o\Bigl(\frac{\ln n}{n}\Bigr).
\]
(Vérification en $n = 100$~: la racine est $x \approx 95.4415$~; la
formule à trois termes donne $100 - 4.6052 + 0.0461 = 95.4409$, celle
à deux termes $95.3948$ --- chaque passe gagne l'\omterm{def:b2:structures:generated}{ordre} prévu.) Même
boucle, troisième paysage~: la méthode du~%
\cref{met:b2:comparison:implicit} ne se soucie pas de l'allure de
l'équation, seulement de ce que chaque passe isole l'inconnue
dominante.
\end{example}

\section{Exercices}

\begin{exercise}[$\star$]\label{exo:b2:comparison:1}
Développer en $+\infty$, deux termes au-delà du terme principal~:
\[
\sqrt{x^2 + x + 1} ,
\qquad
\ln(x^2 + x) - 2\ln x,
\qquad
\frac{x + \sin x}{x - \ln x} .
\]
\end{exercise}

\begin{exercise}[$\star$]\label{exo:b2:comparison:2}
Donner la nature (convergence/divergence) et, en cas de divergence,
l'asymptotique dominante de $\sum_{k \leq n} k^\alpha$ pour $\alpha >
-1$, $\alpha = -1$, $\alpha < -1$, via le~%
\cref{thm:b2:comparison:seriesintegral}.
\end{exercise}

\begin{exercise}[$\star\star$]\label{exo:b2:comparison:3}
Démontrer $H_n = \ln n + \gamma + \frac{1}{2n} + o\bigl(\frac1n\bigr)$.
\emph{(Étudier $v_n = H_n - \ln n - \gamma$~: montrer $v_n - v_{n+1} =
\frac{1}{2n^2} + O(n^{-3})$ et sommer le reste, en comparant avec
$\sum_{k \geq n} \frac{1}{2k^2} \sim \frac{1}{2n}$ ---
\cref{thm:b2:comparison:seriesintegral}~(1).)}
\end{exercise}

\begin{exercise}[$\star\star$]\label{exo:b2:comparison:4}
À l'aide de Stirling, trouver des équivalents de~: $\dfrac{(3n)!}{(n!)^3}$;
$\;\dfrac{n!}{n^n}$; $\;\sqrt[n]{n!}$ (sous la forme $\frac n\eu(1 + o(1))$,
précisée à deux termes).
\end{exercise}

\begin{exercise}[$\star\star$]\label{exo:b2:comparison:5}
Pour $n \geq 2$, démontrer que $x^n + x = 1$ a une unique solution $x_n
\in \intoo{0}{1}$, que $x_n \to 1$, et établir
\[
x_n = 1 - \frac{\ln n}{n} + o\Bigl(\frac{\ln n}{n}\Bigr).
\]
\emph{(De $x_n^n = 1 - x_n$~: prendre les logarithmes et amorcer avec
$x_n = 1 - \varepsilon_n$.)}
\end{exercise}

\begin{exercise}[$\star\star$]\label{exo:b2:comparison:6}
Pousser l'\cref{ex:b2:comparison:tan} d'un \omterm{def:b2:structures:generated}{ordre} de plus~:
\[
x_n = n\pi + \frac\pi2 - \frac{1}{n\pi} + \frac{1}{2n^2\pi} +
o\Bigl(\frac{1}{n^2}\Bigr).
\]
\emph{(Écrire $\cot y_n = x_n$ exactement, développer $\cot y = \frac1y -
\frac y3 + o(y)$ et $x_n = n\pi(1 + \frac{1}{2n} - \dots)$, et
identifier.)}
\end{exercise}

\begin{exercise}[$\star\star$]\label{exo:b2:comparison:7}
Déterminer $\lim_{n\to\infty} \dfrac{1}{n!}\sum_{k=0}^{n} k!$
\emph{(majorer la somme de tous les termes sauf les deux derniers)}, et
en déduire le \omterm{def:b2:comparison:expansion}{développement asymptotique} $\sum_{k \leq n} k! = n!\bigl(1 + \frac1n +
O(n^{-2})\bigr)$.
\end{exercise}

\begin{exercise}[$\star\star\star$]\label{exo:b2:comparison:8}
Soit $u_0 > 0$ et $u_{n+1} = u_n + \dfrac{1}{u_n}$. Démontrer que $u_n
\to \infty$, puis que $u_n \sim \sqrt{2n}$ \emph{(étudier $u_n^2$~:
ses accroissements sont $2 + u_n^{-2}$~; sommer)}, et raffiner~:
\[
u_n = \sqrt{2n}\Bigl(1 + \frac{\ln n}{8n} +
o\Bigl(\frac{\ln n}{n}\Bigr)\Bigr).
\]
\emph{(De $u_n^2 = 2n + \sum_{k<n} u_k^{-2} + u_0^2$ et $u_k^2
\sim 2k$~: la somme est $\sim \frac12\ln n$ d'après le~%
\cref{thm:b2:comparison:seriesintegral}.)}
\end{exercise}

\begin{exercise}[$\star\star\star$]\label{exo:b2:comparison:9}
(Une somme de Riemann avec une astuce) Déterminer le comportement
asymptotique de
\[
S_n = \sum_{k=1}^{n} \frac{1}{n + k\ln n} .
\]
\emph{(Factoriser $n$~: $S_n = \frac1n\sum_k \bigl(1 +
\frac{k\ln n}{n}\bigr)^{-1}$~; reconnaître une somme de type Riemann
avec un paramètre lentement variable $t = \ln n$, calculer $\int_0^1
\frac{\dd u}{1 + tu} = \frac{\ln(1+t)}{t}$, et conclure $S_n \sim
\frac{\ln\ln n}{\ln n}$.)}
\end{exercise}

\begin{exercise}[$\star$]\label{exo:b2:comparison:10}
Démontrer l'identité $(\ln n)^{\ln n} = n^{\ln\ln n}$, puis classer les
suivants par \omterm{def:b2:structures:generated}{ordre} $o(\cdot)$ croissant à l'infini, avec
démonstrations~: $n^2$, $(\ln n)^{\ln n}$, $2^n$, $n!$, $n^n$.
\end{exercise}

\begin{exercise}[$\star\star$]\label{exo:b2:comparison:11}
(Reste de $\sum 1/k^2$, deux termes) À l'aide du télescopage exact
$\sum_{k > n} \frac{1}{k(k+1)} = \frac{1}{n+1}$ et de la
décomposition $\frac1{k^2} = \frac{1}{k(k+1)} +
\frac{1}{k^2(k+1)}$, démontrer
\[
\sum_{k > n} \frac{1}{k^2}
= \frac1n - \frac{1}{2n^2} + O\Bigl(\frac{1}{n^3}\Bigr).
\]
\end{exercise}

\begin{exercise}[$\star\star\star$]\label{exo:b2:comparison:12}
Soit $u_0 = \frac12$ et $u_{n+1} = u_n + \eu^{-u_n}$. Démontrer que
$u_n \to \infty$, puis --- en posant $v_n = \eu^{u_n}$ et en montrant
$v_{n+1} = v_n + 1 + \frac{1}{2v_n} + O\bigl(v_n^{-2}\bigr)$ ---
établir
\[
u_n = \ln n + \frac{\ln n}{2n} + O\Bigl(\frac1n\Bigr).
\]
\end{exercise}

\section{Problème~: amorçage, d'Euler--Maclaurin aux nombres
premiers}

Une quantité implicite ou accumulée livre rarement son
asymptotique d'un coup~; on l'extrait par passes, chaque passe
réinjectant l'estimation précédente dans la relation qui la définit.
Ce problème du week-end entraîne cette boucle sur des équations
nouvelles, démontre la \emph{formule d'Euler--Maclaurin} au premier
\omterm{def:b2:structures:generated}{ordre} (la version trapèze de la comparaison série--intégrale, avec des
barres d'erreur rigoureuses), inverse $x\ln x = n$, et encaisse le
chèque le plus célèbre de la méthode~: à partir du théorème des
nombres premiers admis, la loi asymptotique $p_n \sim n\ln n$ du
$n$-ième nombre premier.

\begin{problem}[{Problème du week-end --- la correction d'Euler--Maclaurin
et l'asymptotique du $n$-ième nombre premier}]
\label{pb:b2:comparison:1}

\medskip
\noindent\textbf{Partie I --- La boucle d'amorçage sur une équation
nouvelle.}
\begin{enumerate}
  \item Démontrer l'affirmation d'unicité de la~%
        \cref{def:b2:comparison:expansion}~: si $f = \sum_{i\leq
        k} c_i\varphi_i + o(\varphi_k) = \sum_{i \leq k}
        c_i'\varphi_i + o(\varphi_k)$ le long de la même échelle, alors
        $c_i = c_i'$ pour tout $i$. Puis pousser l'exemple mixte du
        cours d'un échelon de plus~:
        \[
        \frac{1}{x - \ln x} = \frac1x + \frac{\ln x}{x^2} +
        \frac{(\ln x)^2}{x^3} + o\Bigl(\frac{(\ln
        x)^2}{x^3}\Bigr) \qquad (x \to +\infty),
        \]
        et expliquer pourquoi aucun terme $\frac{c}{x^2}$ n'apparaît.
  \item Montrer que pour tout $n \geq 1$ l'équation $\eu^x + x =
        n$ a exactement une solution réelle $x_n$, et que $x_n
        \to +\infty$ avec $x_n \sim \ln n$.
  \item Amorcer deux fois~:
        \[
        x_n = \ln n - \frac{\ln n}{n} - \frac{(\ln n)^2}{2n^2} +
        o\Bigl(\frac{(\ln n)^2}{n^2}\Bigr).
        \]
  \item Vérifier numériquement en $n = 1000$~: comparer $x_{1000}
        \approx 6.90083$ avec les valeurs à un, deux et trois termes
        de la question 3, à cinq décimales.
\end{enumerate}

\medskip
\noindent\textbf{Partie II --- Euler--Maclaurin, \omterm{def:b2:structures:generated}{ordre} un.}
\begin{enumerate}[resume]
  \item Démontrer l'identité du noyau trapèze~: pour $g$ de classe
        $C^2$ sur $\intcc{0}{1}$,
        \[
        \int_0^1 g(t)\,\dd t = \frac{g(0) + g(1)}{2}
        - \frac12\int_0^1 t(1 - t)\,g''(t)\,\dd t
        \]
        \emph{(intégrer $\frac12 t(1-t)g''$ par parties deux fois)}.
  \item Soit $f$ de classe $C^2$ sur $\intco{1}{+\infty}$ avec
        $\int_1^\infty \abs{f''} < \infty$. Montrer que
        \[
        E_n = \sum_{k=1}^{n} f(k) - \int_1^n f -
        \frac{f(1) + f(n)}{2}
        \]
        converge vers une constante $E$, avec la majoration du reste
        $\abs{E - E_n} \leq \frac18\int_n^\infty\abs{f''}$~: la
        \emph{formule d'Euler--Maclaurin} au premier \omterm{def:b2:structures:generated}{ordre}.
  \item Appliquer ceci à $f(t) = \frac1t$~: démontrer
        \[
        H_n = \ln n + \gamma + \frac{1}{2n} + \varepsilon_n,
        \qquad \abs{\varepsilon_n} \leq \frac{1}{8n^2},
        \]
        renforçant l'\cref{exo:b2:comparison:3} (identifier la
        constante à $\gamma$ en comparant avec
        l'\cref{ex:b2:comparison:harmonic}).
  \item Extraire le coefficient suivant~: montrer $\varepsilon_n =
        -\frac{1}{12n^2} + o\bigl(\frac1{n^2}\bigr)$
        \emph{(les accroissements de $E_n$ valent
        $\frac12\int_0^1t(1-t)f''(n+t)\dd t =
        \frac1{12}f''(n) + o(f''(n))$~; sommer le reste avec le~%
        \cref{thm:b2:comparison:seriesintegral})}.
  \item Appliquer la question 6 à $f = \ln$~: redémontrer en trois
        lignes la convergence de $d_n = \ln n! - (n +
        \frac12)\ln n + n$ (étape 1 du~%
        \cref{thm:b2:comparison:stirling}), avec le taux d'erreur en
        prime $d_n = d + O\bigl(\frac1n\bigr)$.
  \item Appliquer la question 6 à $f(t) = \frac{1}{\sqrt t}$~: montrer
        \[
        \sum_{k=1}^{n}\frac1{\sqrt k} = 2\sqrt n + c +
        \frac{1}{2\sqrt n} + O\Bigl(\frac{1}{n^{3/2}}\Bigr)
        \]
        pour une certaine constante $c$, et évaluer tous les termes en
        $n = 10^4$ (la constante est $c \approx -1.4604$).
\end{enumerate}

\medskip
\noindent\textbf{Partie III --- Inversion~: l'équation $x\ln x =
n$.}
\begin{enumerate}[resume]
  \item Montrer que $x\ln x = n$ a exactement une solution $x_n \in
        \intco{1}{+\infty}$ pour $n \geq 1$, que $x_n \to
        \infty$, et que $\ln x_n \sim \ln n$.
  \item En déduire l'inversion à un terme $x_n \sim
        \dfrac{n}{\ln n}$, puis amorcer encore une fois~:
        \[
        \ln x_n = \ln n - \ln\ln n + o(1),
        \qquad
        x_n = \frac{n}{\ln n}\Bigl(1 + \frac{\ln\ln n}{\ln n}
        + o\Bigl(\frac{\ln\ln n}{\ln n}\Bigr)\Bigr).
        \]
  \item Tester en $n = 10^6$~: la vraie racine est $x \approx
        87\,848$~; comparer avec les valeurs à un terme ($\approx
        72\,382$) et à deux termes ($\approx 86\,140$), et expliquer
        la lenteur du gain (le paramètre du développement est
        $\frac{\ln\ln n}{\ln n}$, seulement $\approx 0.19$ en $n =
        10^6$).
  \item Nous \emph{admettons} maintenant le théorème des nombres
        premiers~: le nombre $\pi(x)$ de nombres premiers $\leq x$
        vérifie $\pi(x) \sim \frac{x}{\ln x}$ quand $x \to \infty$
        (démontré honnêtement dans le volume de troisième année). En
        notant $p_n$ le $n$-ième nombre premier, justifier $\pi(p_n) =
        n$, et faire tourner l'inversion des questions 11--12 pour
        démontrer
        \[
        p_n \sim n \ln n .
        \]
  \item Dividendes~: (a) montrer $\sum_{k \leq n} p_k \sim
        \frac{n^2\ln n}{2}$ \emph{(comparer $\sum k\ln k$ avec
        $\int t\ln t\,\dd t$)}~; (b) calculer la probabilité
        approximative qu'un entier tiré uniformément au hasard à $100$
        chiffres soit premier ($\ln 10^{100} \approx 230.26$~: environ
        un sur $230$).
\end{enumerate}

\medskip
\noindent\textbf{Partie IV --- La méthode exportée~: $x\tan x =
1$.}
\begin{enumerate}[resume]
  \item Montrer que pour tout $n \geq 1$ l'équation $\tan x =
        \frac1x$ a exactement une solution $x_n$ dans
        $\intoo{n\pi}{\,n\pi + \frac\pi2}$, et que $z_n = x_n
        - n\pi \to 0^+$.
  \item Un terme~: $z_n \sim \dfrac{1}{n\pi}$.
  \item Montrer que le développement de $z_n$ n'a \emph{aucun}
        terme $\frac{c}{n^2}$~: $z_n = \frac1{n\pi} +
        O\bigl(\frac{1}{n^3}\bigr)$.
  \item Trois termes~: à l'aide de $\arctan u = u - \frac{u^3}3 +
        O(u^5)$ et $\frac1{x_n} = \frac{1}{n\pi} -
        \frac{z_n}{(n\pi)^2} + O(n^{-3}\cdot z_n^2)$, démontrer
        \[
        x_n = n\pi + \frac{1}{n\pi} -
        \frac{4}{3\pi^3 n^3} + o\Bigl(\frac{1}{n^3}\Bigr).
        \]
  \item Vérifier en $n = 3$~: vraie racine $x_3 \approx 9.5293344$~;
        comparer les valeurs à un et à trois termes, et contraster en
        une phrase avec le $\tan x = x$ du cours
        (\cref{ex:b2:comparison:tan})~: où chaque suite se place
        dans sa fenêtre, et pourquoi.
\end{enumerate}

\medskip
\noindent\textbf{Partie V --- Un amorçage dynamique, règles du
jeu, synthèse.}
\begin{enumerate}[resume]
  \item Soit $u_0 \in \intoo{0}{\pi}$ et $u_{n+1} = \sin u_n$.
        Montrer que $u_n \to 0$ en décroissant, et calculer la limite
        de $\dfrac{1}{u_{n+1}^2} - \dfrac{1}{u_n^2}$ \emph{(développer
        $\sin^{-2}$ via $\sin u = u - \frac{u^3}6 + o(u^3)$)}.
  \item En déduire, via les moyennes de Cesàro (volume de première
        année), le classique
        \[
        u_n \sim \sqrt{\frac{3}{n}} .
        \]
  \item (Numérique certifié) À l'aide de la majoration rigoureuse de
        la question 7, montrer qu'évaluer $\ln n + \gamma + \frac1{2n}$
        en $n = 10^6$ donne $H_{10^6}$ avec une erreur d'au plus
        $1.25\cdot10^{-13}$ --- une somme d'un million de termes
        calculée à treize chiffres par trois termes.
  \item (Règles du jeu) Démontrer ou réfuter, avec démonstrations ou
        contre-exemples~: (a) si $u_n \sim v_n \to +\infty$ alors
        $\ln u_n \sim \ln v_n$~; (b) si $u_n \sim v_n$ alors
        $\eu^{u_n} \sim \eu^{v_n}$~; (c) si $f \sim g$ en
        $+\infty$ ($f, g$ dérivables) alors $f' \sim g'$.
  \item (Synthèse) En une phrase chacun~: la boucle d'amorçage du~%
        \cref{met:b2:comparison:implicit} telle qu'utilisée dans les
        parties I, III, IV~; ce que la correction trapèze ajoute au~%
        \cref{thm:b2:comparison:seriesintegral}~; pourquoi l'inversion
        de $x\ln x$ est exactement le pont de $\pi(x)$ à
        $p_n$~; et laquelle des règles de la question 24 a protégé
        quelle étape. Nommer les deux sommets~: la formule
        d'Euler--Maclaurin (au premier \omterm{def:b2:structures:generated}{ordre}), et la loi asymptotique
        du $n$-ième nombre premier.
\end{enumerate}
\end{problem}
